\begin{abstract}
Algorithms that favor popular items are used to help us select among many choices, from engaging articles on a social media news feed to songs and books that others have purchased, and from top-raked search engine results to highly-cited scientific papers. The goal of these algorithms is to identify high-quality items such as reliable news, beautiful movies, prestigious information sources, and important discoveries --- in short, high-quality content should rank at the top.  
Prior work has shown that choosing what is popular may amplify random fluctuations and ultimately lead to sub-optimal rankings. Nonetheless, it is often assumed that recommending what is popular will help high-quality content ``bubble up'' in practice. Here we identify the conditions in which popularity may be a viable proxy for quality content by studying a simple model of cultural market endowed with an intrinsic notion of quality. A parameter representing the cognitive cost of exploration controls the critical trade-off between quality and popularity.  We find a regime of intermediate exploration cost where an optimal balance exists, such that choosing what is popular actually promotes high-quality items to the top. Outside of these limits, however, popularity bias is more likely to hinder quality. 
These findings clarify the effects of algorithmic popularity bias on quality outcomes, and may inform the design of more principled mechanisms for techno-social cultural markets.
\end{abstract}